\documentclass[10pt,a4paper]{amsart}
\usepackage[margin=19mm]{geometry}
\usepackage{amsmath,amssymb,mathtools}
\usepackage[scr=rsfs,cal=euler]{mathalfa}
\usepackage{xfrac}
\usepackage[unicode,hidelinks,bookmarks=false]{hyperref}
\newcommand{\ie}{i.e.\ }
\newcommand{\cf}{cf.\ }
\newcommand{\resp}{respectively\ }
\newcommand{\Subsection}{\subsection}
\providecommand{\nobreakdash}{\nobreak\hbox{-}}
\setlength{\emergencystretch}{2em}
\multlinegap0pt
\usepackage{xcolor}
\usepackage{algorithm}\usepackage{algpseudocode}
\usepackage{amssymb}
\usepackage{dsfont}
\newtheorem{assumption}{Assumption}
\newtheorem{proposition}{Proposition}
\title{Flow Matching for Measure Transport and Feedback Control of Control-Affine Systems}
\author{Karthik Elamvazhuthi}
\date{Source: arXiv:2510.02706v3 (CC BY 4.0)}
\begin{document}
\maketitle

\noindent\emph{Source and licence.} Flow Matching for Measure Transport and Feedback Control of Control-Affine Systems. Version arXiv:2510.02706v3 (CC BY 4.0), distributed by arXiv under the Creative Commons Attribution 4.0 International licence, http://creativecommons.org/licenses/by/4.0/, as recorded in the arXiv licence metadata for this exact version and shown on the abstract page https://arxiv.org/abs/2510.02706v3. Full licence notice: \texttt{licenses/arxiv-2510.02706-CC-BY-4.0.txt}.\par\smallskip

\noindent\textit{Editorial scope.} Counted unit: the article's proposition prop:relaass (source0.tex lines 2449--2454). The article's complete original proof of it is delivered with it (source0.tex lines 2456--2559, one \texttt{proof} environment of 104 lines). No mathematics is written, repaired or re-derived: the statement and the proof are the article's own lines, in the article's own notation, and no retained line is substituted or re-flowed.  Retained uncounted as the source-local context the proof needs: lines 1627--1634; lines 1639--1648.  Nothing was omitted from any retained span, so the package carries no omission record.  No retained line contains a citation, so the wrapper carries no bibliography and no bibliography entry.  No retained line contains a figure environment, an image-inclusion command or drawing code, so no image is delivered and the package declares no asset dependency.  Editorial formatting only, in addition: an AMS \texttt{amsart} preamble that restates the article's own declarations and macros used by the retained lines, namely the environments \texttt{assumption}, \texttt{proposition} on separate counters, together with the standard packages those lines need.  The retained environments print in this document as Assumption 1, Proposition 1 in order of appearance, whereas the article prints the same items with the numbering of its own full document.  The complete native source of the article as distributed in the arXiv e-print for this exact version is bundled unchanged as \texttt{sources/186304/source0.tex}.
\begin{assumption}
\label{asmp:pmp-control}
The map $\mathbb{R}^d \times \mathbb{R}^d \ni (\omega, p) \mapsto \alpha(\omega,p) \in \mathbb{R}^m$ defined by
\[\alpha(\omega,p) := \arg\min_{\alpha \in \mathbb{R}^m} \left[\langle p, f(\omega,\alpha) \rangle + L(\omega,\alpha)\right]\]
is uniquely defined and locally Lipschitz continuous. Moreover, there exists a constant $M > 0$ such that 
\[|\alpha(\omega,p)| \leq M(1 + |\omega| + |p|)\]
for all $(\omega, p) \in \mathbb{R}^d \times \mathbb{R}^d$.
\end{assumption}

\begin{assumption}
\label{asmp:cost}
The cost function $L: \mathbb{R}^d \times \mathbb{R}^m \rightarrow \mathbb{R}$ satisfies:
\begin{enumerate}
\item There exists $\theta > 0$ such that $L(x,u) \geq \theta |u|^2$ for all $(x,u) \in \mathbb{R}^d \times \mathbb{R}^m$.
\item $L \in C^2(\mathbb{R}^d \times \mathbb{R}^m)$, and $\nabla_x L(x, u)$ has uniform linear growth in $x$: there exists $C > 0$ such that
\[|\nabla_x L(x, u)| \leq C(1 + |x|)\]
for all $(x, u) \in \mathbb{R}^d \times \mathbb{R}^m$.
\end{enumerate}
\end{assumption}

\begin{proposition}
    Suppose $f_i :\mathbb{R}^d \rightarrow \mathbb{R}^d$ are smooth, globally bounded vector-fields for $i =1,...,m$, and the cost function $L: \mathbb{R}^d \times \mathbb{R}^m \rightarrow \mathbb{R}$ is such that $L(x, \cdot )$ is strongly convex, uniformly in $x$. Additionally, suppose Assumption \ref{asmp:cost} holds. Then, consider the function,
    \[\alpha(\omega,p) := \arg\min_{\alpha \in \mathbb{R}^m} \left[\langle p, f(\omega,\alpha) \rangle + L(\omega,\alpha)\right]\]
    Then Assumption \ref{asmp:pmp-control} holds true.
    \label{prop:relaass}
\end{proposition}

\begin{proof}

Since $L(\omega,\cdot)$ is uniformly strongly convex and $
a \mapsto \sum_{i=1}^m a_i \langle p,f_i(\omega)\rangle $
is affine, the function
$
a \mapsto \Phi(\omega,p,a): =\langle p,f(\omega,a)\rangle + L(\omega,a)
$
is strongly convex for every $(\omega,p)$. Moreover, because the vector fields $f_i$ are globally bounded and Assumption \ref{asmp:cost} gives quadratic growth in the control variable, $\Phi(\omega,p,\cdot)$ is coercive. Hence $\Phi(\omega,p,\cdot)$ admits a unique minimizer, which we denote by $\alpha(\omega,p)$.

We show local Lipschitz continuity using a standard implicit function theorem argument.
Define
$G(\omega,p,a):=\nabla_a \Phi(\omega,p,a).$
Since $f$ is affine in $a$,
\[
G(\omega,p,a)
=
\bigl(\langle p,f_1(\omega)\rangle,\dots,\langle p,f_m(\omega)\rangle\bigr)
+
\nabla_a L(\omega,a).
\]
The map $G$ is $C^1$ because the vector fields are smooth and $L\in C^2$.

By construction, $\alpha(\omega,p)$ satisfies the first-order optimality condition
\[
G(\omega,p,\alpha(\omega,p))=0.
\]
Moreover,
\[
D_aG(\omega,p,a)=D^2_{aa}L(\omega,a).
\]
Since $L(\omega,\cdot)$ is strongly convex uniformly in $\omega$, there exists $\lambda>0$ such that
\[
\xi^\top D^2_{aa}L(\omega,a)\,\xi \ge \lambda |\xi|^2
\qquad
\text{for all }(\omega,a,\xi)\in\mathbb R^d\times\mathbb R^m\times\mathbb R^m.
\]
Hence $D_aG(\omega,p,a)$ is invertible for every $(\omega,p,a)$.
Applying the implicit function theorem at each point $
(\omega,p,\alpha(\omega,p))$, we conclude that $(\omega,p)\mapsto \alpha(\omega,p)$ is $C^1$ locally.

Next, we derive the linear growth estimate.
By minimality of $\alpha(\omega,p)$,
\[
\Phi(\omega,p,\alpha(\omega,p))\le \Phi(\omega,p,0).
\]
This becomes equivalent to,
\[
L(\omega,\alpha(\omega,p))
+
\sum_{i=1}^m \alpha_i(\omega,p)\,\langle p,f_i(\omega)\rangle
\le
L(\omega,0).
\]
Using again Assumption \ref{asmp:cost} and boundedness of the $f_i$,
\[
\theta |\alpha(\omega,p)|^2
\le
L(\omega,0)+ C_f\sqrt m\,|p|\,|\alpha(\omega,p)|.
\]
with $C_f$ the maximum of the upper bounds of $\sup_{x \in \mathbb{R}^d} |f_i(x)|$.
We will now bound $L(\omega,0)$. Since Assumption \ref{asmp:cost} gives
\[
|\nabla_xL(x,u)|\le C(1+|x|)
\qquad \text{for all }(x,u) \in \mathbb{R}^d \times \mathbb{R}^m,
\]
we may integrate along the segment $s\mapsto s\omega$ to obtain
\[
|L(\omega,0)-L(0,0)|
\le
\int_0^1 |\nabla_xL(s\omega,0)\cdot \omega|\,ds
\le
C|\omega|\int_0^1 (1+s|\omega|)\,ds
\le
C(1+|\omega|^2).
\]
Therefore
\[
L(\omega,0)\le C(1+|\omega|^2).
\]
Substituting this into the previous estimate yields
\[
\theta |\alpha(\omega,p)|^2
\le
C(1+|\omega|^2)+C|p|\,|\alpha(\omega,p)|.
\]
Applying Young's inequality to the last term,
\[
C|p|\,|\alpha(\omega,p)|
\le
\frac{\theta}{2}|\alpha(\omega,p)|^2 + C|p|^2,
\]
and hence
\[
|\alpha(\omega,p)|^2 \le C(1+|\omega|^2+|p|^2).
\]
Taking square roots,
\[
|\alpha(\omega,p)|\le M(1+|\omega|+|p|),
\]
for some constant $M>0$.

This proves Assumption \ref{asmp:pmp-control} holds.
\end{proof}

\end{document}
